\documentclass[10pt, letterpaper]{article}

% Packages:
\usepackage[
    ignoreheadfoot, % set margins without considering header and footer
    top=2 cm, % seperation between body and page edge from the top
    bottom=2 cm, % seperation between body and page edge from the bottom
    left=2 cm, % seperation between body and page edge from the left
    right=2 cm, % seperation between body and page edge from the right
    footskip=1.0 cm, % seperation between body and footer
    % showframe % for debugging
]{geometry} % for adjusting page geometry
\usepackage[explicit]{titlesec} % for customizing section titles
\usepackage{tabularx} % for making tables with fixed width columns
\usepackage{array} % tabularx requires this
\usepackage[dvipsnames]{xcolor} % for coloring text
\definecolor{primaryColor}{RGB}{0, 79, 144} % define primary color
\usepackage{enumitem} % for customizing lists
\usepackage{fontawesome5} % for using icons
\usepackage{amsmath} % for math
\usepackage[
    pdftitle={Andrés Ortega CV},
    pdfauthor={Andrés González Ortega},
    pdfcreator={LaTeX with RenderCV},
    colorlinks=true,
    urlcolor=primaryColor
]{hyperref} % for links, metadata and bookmarks
\usepackage[pscoord]{eso-pic} % for floating text on the page
\usepackage{calc} % for calculating lengths
\usepackage{bookmark} % for bookmarks
\usepackage{lastpage} % for getting the total number of pages
\usepackage{changepage} % for one column entries (adjustwidth environment)
\usepackage{paracol} % for two and three column entries
\usepackage{ifthen} % for conditional statements
\usepackage{needspace} % for avoiding page brake right after the section title
\usepackage{iftex} % check if engine is pdflatex, xetex or luatex

% Ensure that generate pdf is machine readable/ATS parsable:
\ifPDFTeX
    \input{glyphtounicode}
    \pdfgentounicode=1
    \usepackage[T1]{fontenc}
    \usepackage[utf8]{inputenc}
    \usepackage{lmodern}
\fi

\usepackage[default, type1]{sourcesanspro}

% Some settings:
\AtBeginEnvironment{adjustwidth}{\partopsep0pt} % remove space before adjustwidth environment
\pagestyle{empty} % no header or footer
\setcounter{secnumdepth}{0} % no section numbering
\setlength{\parindent}{0pt} % no indentation
\setlength{\topskip}{0pt} % no top skip
\setlength{\columnsep}{0.15cm} % set column seperation
\makeatletter
\let\ps@customFooterStyle\ps@plain % Copy the plain style to customFooterStyle
\patchcmd{\ps@customFooterStyle}{\thepage}{
    \color{gray}\textit{\small Andrés Ortega - Pagina \thepage{} de \pageref*{LastPage}}
}{}{} % replace number by desired string
\makeatother
\pagestyle{customFooterStyle}

\titleformat{\section}{
    % avoid page braking right after the section title
    \needspace{4\baselineskip}
    % make the font size of the section title large and color it with the primary color
    \Large\color{primaryColor}
}{
}{
}{
    % print bold title, give 0.15 cm space and draw a line of 0.8 pt thickness
    % from the end of the title to the end of the body
    \textbf{#1}\hspace{0.15cm}\titlerule[0.8pt]\hspace{-0.1cm}
}[] % section title formatting

\titlespacing{\section}{
    % left space:
    -1pt
}{
    % top space:
    0.3 cm
}{
    % bottom space:
    0.2 cm
} % section title spacing

% \renewcommand\labelitemi{$\vcenter{\hbox{\small$\bullet$}}$} % custom bullet points
\newenvironment{highlights}{
    \begin{itemize}[
        topsep=0.10 cm,
        parsep=0.10 cm,
        partopsep=0pt,
        itemsep=0pt,
        leftmargin=0.4 cm + 10pt
    ]
}{
    \end{itemize}
} % new environment for highlights

\newenvironment{highlightsforbulletentries}{
    \begin{itemize}[
        topsep=0.10 cm,
        parsep=0.10 cm,
        partopsep=0pt,
        itemsep=0pt,
        leftmargin=10pt
    ]
}{
    \end{itemize}
} % new environment for highlights for bullet entries


\newenvironment{onecolentry}{
    \begin{adjustwidth}{
        0.2 cm + 0.00001 cm
    }{
        0.2 cm + 0.00001 cm
    }
}{
    \end{adjustwidth}
} % new environment for one column entries

\newenvironment{twocolentry}[2][]{
    \onecolentry
    \def\secondColumn{#2}
    \setcolumnwidth{\fill, 4.5 cm}
    \begin{paracol}{2}
}{
    \switchcolumn \raggedleft \secondColumn
    \end{paracol}
    \endonecolentry
} % new environment for two column entries

\newenvironment{threecolentry}[3][]{
    \onecolentry
    \def\thirdColumn{#3}
    \setcolumnwidth{1 cm, \fill, 4.5 cm}
    \begin{paracol}{3}
    {\raggedright #2} \switchcolumn
}{
    \switchcolumn \raggedleft \thirdColumn
    \end{paracol}
    \endonecolentry
} % new environment for three column entries

\newenvironment{header}{
    \setlength{\topsep}{0pt}\par\kern\topsep\centering\color{primaryColor}\linespread{1.5}
}{
    \par\kern\topsep
} % new environment for the header

\newcommand{\placelastupdatedtext}{% \placetextbox{<horizontal pos>}{<vertical pos>}{<stuff>}
  \AddToShipoutPictureFG*{% Add <stuff> to current page foreground
    \put(
        \LenToUnit{\paperwidth-2 cm-0.2 cm+0.05cm},
        \LenToUnit{\paperheight-1.0 cm}
    ){\vtop{{\null}\makebox[0pt][c]{
        \small\color{gray}\textit{Marzo 2026}\hspace{\widthof{Marzo 2026}}
    }}}%
  }%
}%

% save the original href command in a new command:
\let\hrefWithoutArrow\href

% new command for external links:
\renewcommand{\href}[2]{\hrefWithoutArrow{#1}{\ifthenelse{\equal{#2}{}}{ }{#2 }\raisebox{.15ex}{\footnotesize \faExternalLink*}}}


\begin{document}
    \newcommand{\AND}{\unskip
        \cleaders\copy\ANDbox\hskip\wd\ANDbox
        \ignorespaces
    }
    \newsavebox\ANDbox
    \sbox\ANDbox{}

    \placelastupdatedtext
    \begin{header}
        \fontsize{30 pt}{30 pt}
        \textbf{Andrés González Ortega}

        \vspace{0.3 cm}

        \normalsize
        \mbox{{\footnotesize\faMapMarker*}\hspace*{0.13cm}CDMX}%
        \kern 0.25 cm%
        \AND%
        \kern 0.25 cm%
        \mbox{\hrefWithoutArrow{mailto:gonorandres@ciencias.unam.mx}{{\footnotesize\faEnvelope[regular]}\hspace*{0.13cm}gonorandres@ciencias.unam.mx}}%
        \kern 0.25 cm%
        \AND%
        \kern 0.25 cm%
        \mbox{\hrefWithoutArrow{tel:+52 55 48344672}{{\footnotesize\faPhone*}\hspace*{0.13cm}55 48344672}}%
        \AND%
        \kern 0.25 cm%
        \mbox{\hrefWithoutArrow{https://gonorandres.github.io/}{{\footnotesize\faGithub}\hspace*{0.13cm}Portafolio}}%

    \end{header}

    \vspace{0.3 cm - 0.3 cm}


    \section{Sobre mí}

        \begin{onecolentry}
Soy estudiante de último semestre de la \textbf{Licenciatura en Actuaría} en la Facultad de Ciencias, UNAM, especializado en \textbf{análisis de datos}, \textbf{modelación estadística} y \textbf{gestión de riesgo}. Mi experiencia incluye el desarrollo de modelos predictivos de alto rendimiento, análisis cuantitativo en finanzas y procesamiento de grandes volúmenes de datos. Domino herramientas como \textbf{Python}, \textbf{R} y \textbf{SQL} para transformar datos complejos en insights accionables que soporten la \textbf{toma de decisiones estratégicas}. Me motiva resolver problemas del mundo real aplicando métodos cuantitativos rigurosos y mantengo un enfoque de aprendizaje continuo en tecnologías emergentes del análisis de datos.
        \end{onecolentry}


    \section{Habilidades}

\begin{onecolentry}
    \begin{highlightsforbulletentries}
        \item \textit{Lenguajes}: \textbf{Python}, \textbf{R}, \textbf{C}
        \item \textit{Bases de Datos}: Manejo de \textbf{SQL}
        \item \textit{Herramientas}: \textbf{Excel} (Nivel Intermedio), Microsoft Office
        \item \textit{Análisis y Modelación}: \textbf{Estadística}, \textbf{Modelación de Datos}, Generación de \textbf{Reportes Profesionales}, \textbf{Limpieza de Bases de Datos}
        \item \textit{Formación Técnica}: \textbf{Métodos Cuantitativos en Finanzas}, \textbf{Procesos Estocásticos}, \textbf{Inferencia Estadística}, \textbf{Análisis Multivariado}, Análisis Numérico, Cálculo Aplicado a Optimización, Teoría del Riesgo
        \item \textit{Idiomas}: \textbf{Inglés} (Lectura y Redacción Fluida)
        \item \textbf{Habilidades Personales}: Pensamiento analítico, aprendizaje Autónomo, adaptabilidad, alta capacidad de concentración, gusto por el aprendizaje.
    \end{highlightsforbulletentries}
\end{onecolentry}

    \section{Educación}

        \begin{threecolentry}{\textbf{}}{
            Sept 2021 – En curso
        }
            \textbf{Facultad de Ciencias UNAM}, \textbf{Actuaría}
            \begin{highlights}
                \item Créditos: 87\%
            \end{highlights}
        \end{threecolentry}

    \section{Cursos y certificados}

    \begin{onecolentry}
        \begin{highlightsforbulletentries}
            \item \textbf{Associate Data Analyst/ DataCamp} \href{https://drive.google.com/file/d/1EqzyjRkWOJBlU9zUTpkcUcj6QtuXRZfK/view?usp=sharing}{}
            \item \textbf{Associate Data Scientist R / DataCamp} \href{https://drive.google.com/file/d/1xVLk15A1LBbyXH1fZnd4eBTaJ4d8FLLO/view?usp=sharing}{}
            \item \textbf{Data Scientist in R / Datacamp }\href{https://drive.google.com/file/d/1Woe6xqxofloFZ9uM2f5VsdGxY-WQWCRI/view?usp=sharing}{} \\
            \item \textbf{Examen P (Probabilidad) / SoA} -- Aprobado, Mar 2026 \href{https://drive.google.com/file/d/1Qp8ymh9TK8Y52AvmlXe_pZif3NMRZceA/view?usp=drivesdk}{}
        \end{highlightsforbulletentries}

            Además estoy en proceso de aprendizaje para acreditar lo siguiente:

        \begin{highlightsforbulletentries}
            \item Certificación TOEFL
            \item Data Analyst in Power BI / DataCamp
            \item Examen FM / SoA
        \end{highlightsforbulletentries}
    \end{onecolentry}

\newpage

    \section{Proyectos}

        He realizado diversos proyectos usando el conocimiento adquirido y asumiendo nuevos retos, con el enfoque práctico de controlar la incertidumbre y analizar los posibles riesgos emergentes. Usando técnicas de simulación, estadística y matemáticas en general; destaco algunos de mis favoritos:

\vspace{.2cm}

\begin{twocolentry}{

\href{https://github.com/GonorAndres/Proyectos_Aprendizaje/tree/main/Credit_Risk_Model}{GitHub}

\href{https://drive.google.com/file/d/1vouH3yRX7TZtE7-4qC4kTULfQgg8aQGx/view?usp=sharing}{Resumen Ejectutivo PDF}

\href{https://drive.google.com/file/d/1PrDvAmmp_u9IBGZnbSPdPh7uDb3i5jY5/view?usp=sharing}{PDF Executive Report}
}

}
    \textbf{Modelo Predictivo de Incumplimiento Crediticio}
    \begin{highlights}
        \item Desarrollo de modelo de machine learning con 85.19\% de precisión (AUC = 0.8519) para predecir incumplimiento en préstamos usando \textbf{Python} y técnicas de \textbf{GLM}.
        \item Procesamiento y limpieza de dataset con +32,000 registros, implementando controles de calidad y validación cruzada para asegurar robustez del modelo.
        \item Documentación completa en inglés con enfoque en transparencia regulatoria y aplicabilidad en decisiones crediticias reales.
    \end{highlights}
\end{twocolentry}

\vspace{0.2 cm}

\begin{twocolentry}{

\href{https://github.com/GonorAndres/Proyectos_Aprendizaje/tree/main/Bayesian_vs_Frequentist}{GitHub}

\href{https://drive.google.com/file/d/1i-9ZWjJieLJnb3CnEEdb4FKQd_cIKFff/view?usp=drive_link}{Resumen PDF}

\href{https://drive.google.com/file/d/1kguG7XKBYd6OLJ3nI1z-Syt7Ozcq0mUB/view?usp=drive_link}{PDF Summary}
}

    \textbf{A/B Testing: Inferencia Bayesiana vs Frecuentista}
    \begin{highlights}
        \item Comparación visual e interpretativa entre métodos bayesianos y frecuentistas para estimar tasas de conversión en un test A/B simulado en \textbf{Python}.
        \item Implementación de teoría de decisión con utilidad esperada para justificar decisiones bajo incertidumbre más allá del p-valor clásico.
        \item Documentación completa en inglés, resumen explicativo en LaTeX con lenguaje claro y referencias gráficas.
    \end{highlights}
\end{twocolentry}

\vspace{0.2 cm}

\begin{twocolentry}{

    \href{https://drive.google.com/file/d/1u2uN4X-W-sbW3lEXOigsWpWmMOgzthxX/view?usp=sharing}{Reporte PDF}

}

    \textbf{Limpieza y Transformación de Base de Deuda Pública CDMX}
    \begin{highlights}
        \item Proceso completo de depuración, verificación temporal y estandarización de variables en una base administrativa de deuda pública.
        \item Uso intensivo de funciones de \textbf{Microsoft Excel} como \texttt{BUSCARV()}, \texttt{SI()} y formato condicional para reconstrucción de tasas y fechas.
    \end{highlights}
\end{twocolentry}

\vspace{0.2 cm}

\begin{twocolentry}{
    \href{https://github.com/GonorAndres/Proyectos_Aprendizaje/tree/main/SeriesdeTiempo_TemperaturaDelhi}{GitHub}

    \href{https://drive.google.com/file/d/1-t7sihOVxHWceLPPOFnZBDrdhkBsXKmP/view?usp=drive_link}{Reporte PDF}
}
    \textbf{Series de Tiempo: Temperatura en Delhi (2013-2017)}
    \begin{highlights}
        \item Análisis exploratorio y modelado ARIMA de series de temperatura diaria e identificación de patrones estacionales y tendencias en el clima de Delhi usando \textbf{Software R}.
    \end{highlights}
\end{twocolentry}

\vspace{0.2 cm}

\begin{twocolentry}{
    \href{https://github.com/GonorAndres/Proyectos_Aprendizaje/tree/main/TexasPokerCaseStudy}{GitHub}
}
    \textbf{Estudio de Probabilidades en Texas Hold'em Poker}
    \begin{highlights}
        \item Simulación Monte Carlo para evaluar probabilidades de mano en Texas Hold'em en \textbf{Python}.
        \item Análisis estratégico del juego basado en resultados probabilísticos.
    \end{highlights}
\end{twocolentry}

\vspace{0.4 cm}

\begin{twocolentry}{
    \href{https://github.com/GonorAndres/Proyectos_Aprendizaje/tree/main/EvaluaciónDerivadosDivisas}{GitHub}
}
    \textbf{Evaluación de Productos Derivados de Divisas}
    \begin{highlights}
        \item Construcción de curvas forward y superficies de volatilidad en el mercado FX con los cálculos de precios de forwards y opciones sobre divisas utilizando tablas de volatilidad e interpolación.
    \end{highlights}
\end{twocolentry}

\vspace{0.2 cm}

\begin{twocolentry}{
    \href{https://drive.google.com/file/d/1yi3ml6M0AqRijSQ4R36KaId4pP3GhpME/view?usp=sharing}{Drive Proyecto}

    \href{https://drive.google.com/file/d/1yi3ml6M0AqRijSQ4R36KaId4pP3GhpME/view?usp=sharing}{Reporte PDF}
}
    \textbf{Optimización de Portafolio de Inversión: Teoría de Markowitz}
    \begin{highlights}
        \item Implementación de \textbf{optimización cuadrática} para minimización de riesgo con restricciones de diversificación mínima del 5\% por activo usando \textbf{Excel} y métodos cuantitativos.
        \item Análisis de correlaciones, cálculo de VaR paramétrico (95\%) y pruebas de estrés incluyendo backtesting durante períodos de alta volatilidad.
        \item Gestión de riesgo con monitoreo diario de métricas, superando consistentemente la tasa libre de riesgo (CETES 8.15\%) con control estricto de exposición.
    \end{highlights}
\end{twocolentry}

\vspace{0.2 cm}

\begin{twocolentry}{
    \href{https://www.datacamp.com/datalab/w/6c3df475-55bb-4df0-b5df-78ed52fd99e8/edit}{DataCamp}
}
    \textbf{Análisis de la Edad de Oro de los Videojuegos}
    \begin{highlights}
        \item Extracción y análisis de datos en SQL para identificar tendencias históricas en la industria del videojuego.
    \end{highlights}
\end{twocolentry}

\vspace{0.2 cm}

\begin{twocolentry}{
    \href{https://www.datacamp.com/datalab/w/546fd075-0c3d-4a9c-862a-403e768ce2a4/edit}{DataCamp}
}
    \textbf{Análisis de Deuda Internacional}
    \begin{highlights}
        \item Exploración de indicadores de deuda a nivel global mediante consultas avanzadas en SQL.
    \end{highlights}
\end{twocolentry}

\vspace{0.2 cm}

\section*{Mis repositorios}

\begin{onecolentry}
    \begin{itemize}
    \item \href{https://gonorandres.github.io/}{Mi Página de Proyectos (en constante actualización)}
    \item \href{https://github.com/GonorAndres}{Mis repositorios en GitHub}
    \item \href{https://drive.google.com/drive/folders/1XTLJotSZmg3HoxJYgmhMJrnWhDWz4uqV?usp=sharing}{Drive de mis documentos académicos}
        \item \href{https://drive.google.com/drive/folders/1M7KZaFbTw19vDoA-AeusJDkL--nnAwEN?usp=sharing}{Drive sobre algunos apuntes sueltos}
\end{itemize}
\end{onecolentry}

\end{document}